\section*{CHƯƠNG 3: GIÁ TRỊ, HẠN CHẾ LỊCH SỬ CỦA TẦNG LỚP NHO SĨ THẾ KỶ XVIII VÀ BÀI HỌC ĐỐI VỚI ĐỘI NGŨ TRÍ THỨC VIỆT NAM HIỆN NAY}
\addcontentsline{toc}{section}{CHƯƠNG 3: GIÁ TRỊ, HẠN CHẾ LỊCH SỬ CỦA TẦNG LỚP NHO SĨ THẾ KỶ XVIII VÀ BÀI HỌC ĐỐI VỚI ĐỘI NGŨ TRÍ THỨC VIỆT NAM HIỆN NAY}

\subsection*{3.1. Những đóng góp tiến bộ và hạn chế lịch sử của tầng lớp nho sĩ thế kỷ XVIII dưới góc độ duy vật lịch sử}
\addcontentsline{toc}{subsection}{3.1. Những đóng góp tiến bộ và hạn chế lịch sử của tầng lớp nho sĩ thế kỷ XVIII dưới góc độ duy vật lịch sử}

\subsubsection*{3.1.1. Những giá trị đóng góp về mặt tư tưởng, văn hóa giáo dục và tinh thần độc lập dân tộc}
\addcontentsline{toc}{subsubsection}{3.1.1. Những giá trị đóng góp về mặt tư tưởng, văn hóa giáo dục và tinh thần độc lập dân tộc}

Vận dụng nguyên tắc toàn diện và lịch sử -- cụ thể của chủ nghĩa duy vật lịch sử, có thể khẳng định sự vận động của tầng lớp nho sĩ thế kỷ XVIII đã kết tinh thành bốn giá trị đóng góp nổi bật đối với tiến trình phát triển của dân tộc:

Thứ nhất, trên bình diện tư tưởng triết học và phương pháp luận khoa học, các học giả tiêu biểu thế kỷ XVIII đã bước đầu giải phóng tư duy lý luận Việt Nam khỏi sự ràng buộc của chủ nghĩa duy tâm khách quan và lối học tầm chương trích cú Tống Nho. Việc đề cao bản thể luận vật chất tự phát về ``Khí'', xác lập mệnh đề ``Lý ở trong Khí'' cùng phương pháp khảo chứng thực địa đã đặt nền móng cho lối tư duy trọng thực chứng, hướng tri thức vào giải quyết các vấn đề kinh tế, điển chế, địa lý và xác lập chứng cứ chủ quyền lãnh thổ, biển đảo quốc gia.

Thứ hai, trên bình diện văn hóa -- giáo dục, tầng lớp nho sĩ đã thực hiện thành công quá trình xã hội hóa và bản địa hóa tri thức trong bối cảnh thiết chế giáo dục quan phương suy sụp. Mạng lưới trường tư thục rộng khắp ở nông thôn Bắc Hà và vùng đất mới Gia Định, kết hợp với nỗ lực đưa chữ Nôm vào dịch thuật kinh điển và văn thư hành chính dưới triều Tây Sơn, đã phá vỡ thế độc quyền chữ nghĩa của cung đình, giữ vững truyền thống hiếu học và đạo lý dân tộc ngay từ nền tảng làng xã.

Thứ ba, trên bình diện chính trị và bảo vệ Tổ quốc, bộ phận nho sĩ cấp tiến đã chứng minh vai trò tích cực của ý thức xã hội tiên tiến khi tác động trở lại thực tiễn lịch sử. Bằng việc vượt qua quan niệm trung quân đóng khung trong một dòng họ phong kiến để lấy lợi ích dân tộc và tư tưởng ``dân vi bản'' làm thước đo chính nghĩa, trí tuệ quân sự -- ngoại giao của kẻ sĩ đã kết hợp nhuần nhuyễn với sức mạnh quật khởi của phong trào nông dân, góp phần quyết định vào thắng lợi chống ngoại xâm và giữ vững nền tự chủ quốc gia.

Thứ tư, trên bình diện giá trị nhân văn, sự dung hợp cởi mở giữa Nho -- Phật -- Đạo cùng tiếng nói phê phán lễ giáo phong kiến trong văn học chữ Nôm đã đánh dấu sự thức tỉnh của ý thức cá nhân. Con người hiện thực với nhu cầu hạnh phúc trần thế, đặc biệt là phẩm giá và quyền sống của người phụ nữ, lần đầu tiên trở thành mối quan tâm trung tâm của đời sống tinh thần, tạo bước chuyển quan trọng trong lịch sử văn hóa Việt Nam.

\subsubsection*{3.1.2. Những hạn chế mang tính lịch sử và rào cản thế giới quan giai cấp phong kiến}
\addcontentsline{toc}{subsubsection}{3.1.2. Những hạn chế mang tính lịch sử và rào cản thế giới quan giai cấp phong kiến}

Song song với những đóng góp tiến bộ, do bị quy định bởi tồn tại xã hội đương thời, tầng lớp nho sĩ thế kỷ XVIII vẫn mang những giới hạn lịch sử không thể vượt qua:

Trước hết là giới hạn về thế giới quan chính trị -- giai cấp. Mặc dù kinh tế hàng hóa thế kỷ XVIII đã có bước khởi sắc, xã hội Đại Việt về căn bản vẫn dựa trên phương thức sản xuất phong kiến tiểu nông phân tán, chưa xuất hiện phương thức sản xuất tư bản chủ nghĩa cũng như giai cấp tư sản dân tộc. Bởi vậy, tư duy chính trị của ngay cả những nho sĩ tiến bộ nhất vẫn nằm trong quỹ đạo quân chủ: lý tưởng cao nhất của họ chỉ dừng lại ở mô hình ``minh quân -- nhân chính'' (trông đợi một bậc vua sáng tôi hiền cứu đời), chứ chưa thể hình thành ý niệm về quyền làm chủ của nhân dân hay thiết chế nhà nước dân chủ pháp quyền.

Tiếp đến là giới hạn về trình độ nhận thức khoa học. Trong điều kiện nền sản xuất thủ công và sự vắng mặt của khoa học tự nhiên thực nghiệm hiện đại, những yếu tố duy vật và biện chứng trong học thuyết của các danh nho thế kỷ XVIII mới mang tính chất trực quan, tự phát và kinh nghiệm chủ nghĩa. Khi lý giải các hiện tượng xã hội phức tạp hoặc biến cố tự nhiên, tư duy của họ đôi lúc vẫn pha tạp thuyết thiên mệnh, sấm vĩ và ngũ hành cảm ứng.

Cuối cùng là sức ỳ bảo thủ của ý thức hệ phong kiến lạc hậu. Sự bám rễ sâu của luân lý tam cương và tâm lý đẳng cấp ``kẻ sĩ'' đã khiến một bộ phận không nhỏ trí thức quay lưng lại với phong trào quần chúng: kẻ cố chấp giữ tiết cực đoan đến mức tuẫn vong hoặc tiếp tay cho ngoại bang chống lại dân tộc; người lại thu mình trong lối sống ẩn dật tiêu cực, hoài cổ, gây lãng phí nguồn lực chất xám trước yêu cầu cải cách đất nước.

\subsection*{3.2. Ý nghĩa lịch sử và bài học kinh nghiệm trong việc phát huy vai trò đội ngũ trí thức giai đoạn hiện nay}
\addcontentsline{toc}{subsection}{3.2. Ý nghĩa lịch sử và bài học kinh nghiệm trong việc phát huy vai trò đội ngũ trí thức giai đoạn hiện nay}

\subsubsection*{3.2.1. Bài học về nguyên tắc thống nhất giữa lý luận với thực tiễn và việc rèn luyện bản lĩnh chính trị}
\addcontentsline{toc}{subsubsection}{3.2.1. Bài học về nguyên tắc thống nhất giữa lý luận với thực tiễn và việc rèn luyện bản lĩnh chính trị}

Từ sự thăng trầm của tầng lớp nho sĩ thế kỷ XVIII, bài học đầu tiên có giá trị thời sự trực tiếp là phải quán triệt nghiêm túc nguyên tắc thống nhất giữa lý luận với thực tiễn, kiên quyết loại bỏ bệnh hình thức và tệ sùng bái bằng cấp. Hệ lụy từ việc thương mại hóa tư cách trường thi thời Lê mạt cho thấy: một khi giáo dục tách rời thực tiễn và học vị biến thành công cụ mưu cầu danh lợi thay vì phản ánh năng lực thực chất, nền học thuật sẽ suy thoái từ gốc rễ.

Trong bối cảnh phát triển kinh tế tri thức và chuyển đổi số hiện nay, công tác giáo dục -- đào tạo và quản lý khoa học cần tiếp tục đấu tranh triệt để với bệnh thành tích, tình trạng ``học giả bằng thật'' và thiếu liêm chính học thuật. Kế thừa tinh thần thực học của tiền nhân, giá trị của người trí thức hiện đại phải được kiểm chứng thông qua năng lực sáng tạo công nghệ lõi, các công trình nghiên cứu có khả năng ứng dụng vào sản xuất, quản trị quốc gia và nâng cao chất lượng sống của nhân dân.

Gắn liền với năng lực chuyên môn là bài học về bản lĩnh chính trị và đạo đức trách nhiệm của người trí thức. Thực tiễn thế kỷ XVIII chứng minh rằng trí thức chỉ phát huy được tầm vóc lịch sử khi biết hòa mình vào dòng chảy tiến bộ của dân tộc; ngược lại, nếu đặt lợi ích cục bộ lên trên lợi ích chung hoặc giáo điều, xa rời nhân dân thì tất yếu sẽ rơi vào bế tắc. Trước những biến động đa chiều của tình hình thế giới và mặt trái cơ chế thị trường hiện nay, đội ngũ trí thức Việt Nam cần kiên định mục tiêu độc lập dân tộc gắn liền với chủ nghĩa xã hội trên nền tảng chủ nghĩa Mác -- Lênin, tư tưởng Hồ Chí Minh; nêu cao tinh thần phụng sự Tổ quốc, tích cực tham gia tư vấn, phản biện xã hội khách quan, khoa học nhằm hoàn thiện đường lối phát triển đất nước.

\subsubsection*{3.2.2. Bài học về quan điểm trọng dụng nhân tài và sự vận dụng trong đường lối của Đảng Cộng sản Việt Nam}
\addcontentsline{toc}{subsubsection}{3.2.2. Bài học về quan điểm trọng dụng nhân tài và sự vận dụng trong đường lối của Đảng Cộng sản Việt Nam}

Sự quy tụ đông đảo hiền tài Bắc Hà dưới ngọn cờ Tây Sơn để lại bài học quý giá về nghệ thuật chiêu hiền đãi sĩ trong quản trị quốc gia. Muốn thu hút và phát huy chất xám của đội ngũ trí thức, chủ thể lãnh đạo phải có thái độ cầu thị chân thành, bao dung, không định kiến hẹp hòi về xuất thân hay quá khứ, đồng thời mạnh dạn trao quyền, giao nhiệm vụ xứng tầm để người tài có không gian cống hiến thực chất.

Kế thừa truyền thống ``Hiền tài là nguyên khí của quốc gia'' trên lập trường mácxít, Đảng Cộng sản Việt Nam luôn xác định đội ngũ trí thức là lực lượng lao động sáng tạo đặc biệt quan trọng trong khối liên minh công nhân -- nông dân -- trí thức. Quan điểm nhất quán đó được cụ thể hóa thành tầm nhìn chiến lược tại Nghị quyết số 45-NQ/TW ngày 24/11/2023 của Hội nghị lần thứ tám Ban Chấp hành Trung ương Đảng Khóa XIII về ``Tiếp tục xây dựng và phát huy vai trò của đội ngũ trí thức đáp ứng yêu cầu phát triển đất nước nhanh và bền vững trong giai đoạn mới''.

Nghị quyết số 45-NQ/TW khẳng định xây dựng đội ngũ trí thức vững mạnh toàn diện chính là đầu tư trực tiếp cho bồi đắp nguyên khí quốc gia và phát triển bền vững. Để hiện thực hóa định hướng này, Đảng và Nhà nước tập trung hoàn thiện cơ chế tuyển dụng, trọng dụng, đãi ngộ đặc thù đối với các nhà khoa học đầu ngành và nhân tài trẻ dựa trên hiệu quả cống hiến; kiến tạo môi trường làm việc dân chủ, tôn trọng tự do học thuật gắn với đề cao đạo đức nghề nghiệp; đồng thời xây dựng hành lang pháp lý bảo vệ những cán bộ, trí thức dám nghĩ, dám làm, dám đổi mới sáng tạo vì lợi ích chung. Sự kết hợp hài hòa giữa chính sách đãi ngộ đúng đắn của thể chế với khát vọng dấn thân của bản thân người trí thức chính là động lực quyết định đưa đất nước vươn lên giàu mạnh, văn minh trong kỷ nguyên phát triển mới.
